% Section 19.9, from lectures/L19/README.md.

\section{Summary}
\label{c19:sec:review}

\needspace{6\baselineskip}
\subsection*{What you should be able to do now}
\begin{itemize}
  \item Implement the transaction FSM per the protocol specification, including the abort rule.
  \item Explain why a read latches the register value once, and what goes wrong without that rule.
  \item Compose the whole node into a top level and verify it in simulation.
  \item Wire two DE0-CVs together onto a shared bus, and a microcontroller as master over its MVIO
    domain, and debug the chain systematically from the bottom up.
  \item Say what requirements the SPI connection puts on the hat's design, and why they cannot be
    fixed after the fact.
  \item Say what simulation has proved and what only hardware can prove.
\end{itemize}

\needspace{6\baselineskip}
\subsection*{Questions to test yourself with}
\begin{enumerate}
  \item Why is the register value latched once at the end of the command byte instead of being read
    while the data bytes are shifted out? Which register suffers most from that rule being broken?
  \item SS goes high after three bytes. What has happened to the transaction, and what has changed in
    the node?
  \item A write to a reserved index. What does the bridge do, and why does it consume all five bytes
    anyway?
  \item Why is there no error channel in the transaction layer?
  \item Which bring-up step would reveal a swapped MISO and MOSI, and which would reveal a
    \code{VDDIO2} that was never connected?
  \item The hat puts an LED on \code{PC2}. Which step fails first, and why is that a fault you cannot
    program your way out of?
  \item What did steps 1 to 3 prove that steps 4 to 8 did not have to prove again?
  \item What has the simulation proved about the design, and what has it been unable to prove?
\end{enumerate}

\needspace{6\baselineskip}
\subsection*{Final presentation}
Every group shows, in about fifteen minutes: a run of \code{make build-project} with every testbench
green; a short walkthrough of \code{can_spi_node}, what every block does and where the boundaries
run; the bring-up, as far as you got, on real hardware; and \file{CONTRIBUTORS.md} plus a quick tour
of the repository's history, branches, PRs and reviews.

The final presentation is the project's second and last fixed checkpoint. The project is handed in
the same day; see the checklist in appendix~\ref{app:project}.

The bring-up is not graded. Getting to step 3 of eight is a good result; what is graded is the
simulation behind it.

\needspace{6\baselineskip}
\subsection*{Looking ahead}
\Chapref{c20:ch} is the book's second practical exam, with tasks on the project's modules. After it
the course is over.

The controller you have built is handed over to the parallel class, which writes the C++ driver
against the same registers over the same SPI transport. The shared contract is
appendix~\ref{app:protocol}, and nothing else.
